% ======================================================================
% SECTION 7: AUTONOMOUS QUALITY ASSURANCE AND COMPLIANCE
% File: sections/quality_compliance.tex
% ======================================================================

Quality assurance and compliance provide the foundation for regulatory inspection readiness
and operational integrity across all autonomous sponsor functions. ICH E6(R3) introduced
quality-by-design as a fundamental principle, requiring sponsors to prospectively identify
factors critical to trial quality and implement proportionate quality management systems
\cite{ich-e6r3}. The autonomous quality system operationalizes these requirements through
a dedicated quality agent, a risk-based quality management engine, and electronic systems
governance that ensures 21 CFR Part 11 compliance \cite{cfr-part-11} for all agent-generated
records and decisions.

\subsection{Quality Agent}

The quality agent (quality\_agent) replaces the quality assurance and compliance team. It
automates SOP enforcement, CAPA lifecycle management, audit readiness assessment, change
control workflows, and validation evidence generation. Under ICH E6(R3), the sponsor cannot
delegate accountability for the quality management system even when specific quality
activities are performed by CROs or technology vendors \cite{ich-e6r3}. The quality agent
embodies this non-delegable accountability.

SOP enforcement operates through automated compliance checking. The agent maintains a
repository of all applicable standard operating procedures and continuously verifies that
operational activities conform to procedural requirements. Deviations are detected
through comparison of actual process execution (documented through audit trail entries
from all agents) against the expected process defined in the SOP. When deviations are
detected, they are classified by severity (critical, major, minor) and routed to
appropriate corrective action workflows.

CAPA lifecycle management automates the corrective and preventive action process from
identification through effectiveness verification. Root cause analysis uses structured
methods (five-why analysis, fishbone diagrams, failure mode and effects analysis) to
identify underlying causes of quality events. Corrective actions are assigned with
defined completion timelines and tracked to closure. Effectiveness verification occurs
at a defined interval after corrective action implementation to confirm that the action
resolved the identified issue.

Audit readiness scoring provides a continuous assessment of the organization's preparedness
for regulatory inspection. The agent evaluates readiness across inspection domains: essential
document completeness, audit trail integrity, SOP currency, training documentation, vendor
oversight evidence, safety reporting compliance, and data integrity indicators. A composite
readiness score is computed and reported to the study orchestrator. When the readiness score
drops below the target threshold, the agent generates prioritized remediation recommendations.

\subsection{Risk-Based Quality Management}

Risk-Based Quality Management (RBQM) implements the ICH E6(R3) requirement for
proportionate quality oversight through Quality Tolerance Limits and Key Risk Indicators.
Table~\ref{tab:qtl-kri} presents the QTL/KRI framework for the autonomous sponsor.

\begin{longtable}{L{1.8cm} L{2.5cm} L{2.5cm} L{1.8cm} L{3.0cm}}
\caption{QTL/KRI Framework for Autonomous Sponsor}
\label{tab:qtl-kri} \\
\toprule
\textbf{Quality Domain} & \textbf{QTL Definition} & \textbf{KRI Metric} & \textbf{Threshold} & \textbf{Agent Response} \\
\midrule
\endfirsthead
\toprule
\textbf{Quality Domain} & \textbf{QTL Definition} & \textbf{KRI Metric} & \textbf{Threshold} & \textbf{Agent Response} \\
\midrule
\endhead
\midrule
\multicolumn{5}{r}{\textit{Continued on next page}} \\
\endfoot
\bottomrule
\endlastfoot
Informed consent & All subjects must have valid, current consent before any study procedure & Consent compliance rate by site & $<$ 100\% triggers investigation & Immediate site contact, root cause analysis, re-consent if required \\
SAE reporting timeliness & All SAEs reported within 24 hours of site awareness & Median hours from awareness to initial report & Warning at $>$ 12 hours, action at $>$ 24 hours & Automated escalation to safety agent and sponsor-of-record \\
Protocol deviation rate & Protocol deviations must not exceed acceptable rate & Major deviation rate per 100 subject-visits & Warning at $>$ 3\%, action at $>$ 5\% & Site-level root cause analysis, targeted retraining \\
Data entry timeliness & CRF pages entered within defined window of visit & Median days from visit to CRF completion & Warning at $>$ 5 days, action at $>$ 10 days & Automated query generation, site performance review \\
Query resolution cycle & Data queries resolved within target timeline & Median days from query opening to resolution & Warning at $>$ 7 days, action at $>$ 14 days & Escalation to site, CRO oversight review \\
Monitoring visit completion & Risk-based monitoring visits completed per plan & Percentage of planned visits completed on schedule & Warning at $<$ 90\%, action at $<$ 80\% & Schedule rebalancing, resource reallocation \\
IMP accountability & All IMP units tracked from receipt to disposition & Percentage of IMP with complete chain of custody & $<$ 100\% triggers investigation & Supply agent investigation, site inventory audit \\
\end{longtable}

Centralized statistical monitoring applies statistical methods to accumulating trial data
to detect unusual patterns that may indicate quality problems at specific sites. The quality
agent implements Bayesian hierarchical models that estimate expected data distributions
across sites and flag sites whose data patterns deviate significantly from the cross-site
norm. Monitored patterns include within-site data variability (unusually low variability
may indicate data fabrication), digit preference in laboratory and vital sign measurements,
temporal patterns in data entry (bulk data entry may indicate delayed recording), and
correlation structures between related variables.

\subsection{Electronic Systems Governance}

Electronic systems governance ensures that all autonomous agent operations produce records
compliant with 21 CFR Part 11 \cite{cfr-part-11} and the EMA guideline on computerized
systems used in clinical trials. Every agent decision, data transformation, and communication
generates an electronic record with the required attributes: system-generated timestamp, user
(agent) identification, action description, previous value and new value for data
modifications, and reason for change where applicable.

Table~\ref{tab:electronic-systems} presents the electronic systems validation matrix.

\begin{longtable}{L{2.0cm} L{2.0cm} L{2.0cm} L{3.0cm} L{2.6cm}}
\caption{Electronic Systems Validation Matrix}
\label{tab:electronic-systems} \\
\toprule
\textbf{System} & \textbf{Regulation} & \textbf{Validation Level} & \textbf{Evidence Required} & \textbf{Agent Owner} \\
\midrule
\endfirsthead
\toprule
\textbf{System} & \textbf{Regulation} & \textbf{Validation Level} & \textbf{Evidence Required} & \textbf{Agent Owner} \\
\midrule
\endhead
\midrule
\multicolumn{5}{r}{\textit{Continued on next page}} \\
\endfoot
\bottomrule
\endlastfoot
Agent Decision Engine & 21 CFR Part 11, ICH E6(R3) & Full (GxP Critical) & IQ/OQ/PQ documentation, algorithm validation, audit trail verification & quality\_agent \\
EDC System & 21 CFR Part 11, EMA Annex 11 & Full (GxP Critical) & System validation, data integrity testing, access control verification & data\_biostats\_agent \\
Safety Database & 21 CFR Part 11, ICH E2B(R3) & Full (GxP Critical) & ICSR validation, signal detection algorithm verification, transmission testing & safety\_agent \\
IRT System & 21 CFR Part 11 & Full (GxP Critical) & Randomization validation, supply allocation testing, unblinding controls & supply\_agent \\
Robot Telemetry System & 21 CFR Part 11, adapted 21 CFR 312 & Full (GxP Critical) & Sensor calibration, data integrity, emergency stop validation & robot\_execution\_gateway \\
Document Management & 21 CFR Part 11, ICH E6(R3) & Standard & Access controls, version management, audit trail & writing\_disclosure\_agent \\
Communication System & ICH E6(R3) & Standard & Message integrity, delivery confirmation, retention & study\_orchestrator \\
\end{longtable}

Audit trail integrity verification ensures that the complete decision history of every agent
is preserved in a tamper-evident format. The quality agent periodically verifies audit trail
integrity by computing cryptographic hashes of stored records and comparing them against
reference hashes maintained in the trust layer (\S\ref{sec:trust}). Any detected integrity
violation triggers an immediate investigation with notification to the sponsor-of-record.

Electronic signature management handles the sponsor-of-record approvals required at
mandatory human gates. Each electronic signature complies with 21 CFR Part 11 requirements:
the signature includes the printed name of the signer, the date and time of signing, and the
meaning of the signature (such as review, approval, or responsibility). Electronic
signatures are linked to their respective electronic records such that the records cannot
be altered without invalidating the signature.

The validation lifecycle for each agent system follows a structured approach. Installation
Qualification (IQ) verifies that the agent software is installed correctly in the production
environment with all required dependencies, configurations, and connectivity. Operational
Qualification (OQ) verifies that the agent performs its specified functions correctly across
the range of expected operating conditions, including boundary conditions and error handling.
Performance Qualification (PQ) verifies that the agent performs correctly under
representative production workload conditions with realistic data volumes and concurrent
operations. Each qualification phase produces documented evidence that is maintained in the
quality agent's validation repository.

Periodic revalidation ensures that agent systems remain in a validated state throughout
the trial lifecycle. The quality agent schedules revalidation activities based on a
risk-informed calendar: critical systems (safety database, IRT, robot telemetry)
undergo quarterly revalidation, standard systems (EDC, document management) undergo
semi-annual revalidation, and supporting systems (communication, reporting)
undergo annual revalidation. Change-triggered revalidation occurs whenever an agent
system is updated, reconfigured, or receives a software patch.

The Computer Software Assurance (CSA) approach, replacing the traditional Computer System
Validation (CSV) methodology, provides a risk-based validation framework for the autonomous
sponsor software. Critical thinking and risk assessment determine the level of validation
evidence required for each system component. GxP-critical components (those that directly
affect patient safety, data integrity, or product quality) receive full validation with
scripted testing. Supporting components receive verification through unscripted testing and
operational qualification.
